\section{\evoopd: Lineage-Aware On-Policy Distillation}
\label{sec:method}

\evoopd{} extends on-policy distillation with two additional signals
required for the genome-centric setting: a \emph{verifier-anchored
reward} that exploits the substructure of \genetrace's verifier, and a
\emph{lineage-consistency self-signal} that uses lineage chains to provide
supervision in the absence of gold answers. Both signals are coupled with
the standard teacher reverse-KL via per-token role gating.

\subsection{Per-token role gating}
\label{sec:method:phi}

For each token $t$ in a student rollout $y$, a deterministic
parser~\citep{evo_opd_parser_impl} assigns one of six role tags
$\phi(t) \in \{$boilerplate, content\_field, evidence\_span,
dynamics\_label, gold\_answer, unknown$\}$ based on the rollout's JSON
structure and the GeneTrace schema. Each role has a fixed weight $\alpha(\phi(t))$
(Table~\ref{tab:phi_weights}): boilerplate (JSON brackets, field names)
receives low weight; content fields and dynamics labels receive high
weight. Token-level gating concentrates the optimisation signal on the
tokens that actually carry task semantics.

\begin{table}[t]
  \centering
  \small
  \begin{tabular}{lc}
    \toprule
    Role $\phi$ & Weight $\alpha(\phi)$ \\
    \midrule
    boilerplate     & 0.3 \\
    content\_field  & 1.0 \\
    evidence\_span  & 1.5 \\
    dynamics\_label & 2.0 \\
    gold\_answer    & 0.0 \\
    unknown         & 0.5 \\
    \bottomrule
  \end{tabular}
  \caption{Per-role weights $\alpha(\phi)$. \texttt{gold\_answer} tokens
  are zero-weighted because they are handled separately by the verifier
  head (Sec.~\ref{sec:method:verifier}).}
  \label{tab:phi_weights}
\end{table}

\subsection{Per-token reward}
\label{sec:method:reward}

For each token $t$, the per-token reward combines three signals:

\begin{align}
  r_t = &- \alpha(\phi(t)) \cdot \mathrm{KL}[\pi_\theta \,\|\, \pi_T](y_t) \nonumber \\
        &+ \alpha(\phi(t)) \cdot \lambda_v \cdot v_{\mathrm{adv}}(y, t) \nonumber \\
        &+ \alpha(\phi(t)) \cdot \lambda_c \cdot c_{\mathrm{adv}}(y, p, t)
  \label{eq:rt}
\end{align}

where $\pi_\theta$ is the student, $\pi_T$ is the teacher,
$v_{\mathrm{adv}}$ is the per-token verifier advantage
(Sec.~\ref{sec:method:verifier}), and $c_{\mathrm{adv}}$ is the
per-token lineage-consistency advantage (Sec.~\ref{sec:method:lineage}).
$\lambda_v$ and $\lambda_c$ are scalar loss weights.

The policy is updated by standard PG on
$\sum_t r_t \cdot \nabla \log \pi_\theta(y_t | y_{<t}, x)$ with no critic
and no reference-KL term (these are subsumed by the teacher-KL
contribution).

\subsection{Verifier-anchored reward}
\label{sec:method:verifier}

The verifier $v: y \mapsto [0, 1]$ from \genetrace{}
(Sec.~\ref{sec:genetrace}) is composed of four subscores:
$v_{\mathrm{schema}}$, $v_{\mathrm{evidence}}$, $v_{\mathrm{dynamics}}$,
$v_{\mathrm{match}}$. For closed-form tasks (e.g., \geneexam{}),
$v_{\mathrm{match}}$ is a 0/1 exact match against gold. For open-ended
tasks, $v_{\mathrm{match}}$ is replaced by a continuous PES rubric judge
\citep{ideaevolving2025} that scores Heredity, Variation, and Selection
on a 0--10 scale and normalises to $[0, 1]$.

The per-token advantage distributes the trajectory-level score across
tokens proportional to their role weight:
\begin{equation}
  v_{\mathrm{adv}}(y, t) = (v(y) - \bar{v}) \cdot \mathbb{1}[\phi(t) \in \mathcal{R}_v]
  \label{eq:vadv}
\end{equation}
where $\mathcal{R}_v = \{$content\_field, evidence\_span, dynamics\_label$\}$
and $\bar{v}$ is a running baseline.

\subsection{Lineage-consistency self-signal}
\label{sec:method:lineage}

The lineage-consistency score $c(y, p) \in [0, 1]$ measures whether the
student's output $y$ on paper $p$ is consistent with $p$'s lineage
context. For tasks where $y$ proposes a successor paper, $c$ is computed
by re-prompting the student to predict the \emph{predecessor} of the
proposed successor and comparing to the ground-truth predecessor card
from \genetrace; high $c$ means the proposal is lineage-faithful. For
tasks where $y$ classifies a relationship between two papers, $c$ is
computed by checking the prediction on a held-out third paper in the
same chain.

Critically, $c$ requires no gold answer for $y$ itself --- only the
lineage structure of the \genetrace{} corpus. This makes it applicable
to open-ended tasks where exact-match cannot fire.

\subsection{Open-ended generalisation}
\label{sec:method:open}

The same three-component reward (Eq.~\ref{eq:rt}) applies to open-ended
tasks with two substitutions: $v_{\mathrm{match}}$ becomes the continuous
PES judge, and $c$ uses the backward-prediction variant above. No
algorithmic change is required. We show in
Section~\ref{sec:exp:open_ended} that \evoopd{} improves PES scores on
\geneexam{} \texttt{GENE-Arena} by \todo{X} points over vanilla OPD,
demonstrating that the lineage-consistency signal carries the load when
the verifier alone provides only a noisy rubric score.

When $\lambda_v = \lambda_c = 0$, Eq.~\ref{eq:rt} reduces to vanilla
on-policy distillation with per-token role gating; when additionally
$\alpha \equiv 1$, it reduces to the Tinker recipe \citep{tinker_opd_2025}.
This makes \evoopd{} non-inferior to vanilla OPD by construction (modulo
optimisation noise), and the ablations in Sec.~\ref{sec:exp:ablation}
quantify the contribution of each component.

\subsection{Training data shape}
\label{sec:method:data}

Unlike SFT, \evoopd's training data is regenerated each step from
on-policy rollouts. A single training record contains: the prompt $x$
(sampled from a fixed prompt pool of $\sim$10K, drawn from \genetrace{}
+ math/code preservation + \geneexam{} Arena prompts); the student
rollout $y$; the teacher per-token log-probabilities $\log \pi_T(y_t |
y_{<t}, x)$; the per-token role tags $\phi(t)$ from the parser; the
verifier subscores $v(y)$; the lineage anchor and computed $c(y, p)$;
and the resulting per-token reward $r_t$. Full schema in
Appendix~\ref{app:opd_record}.
